% from ega3-7-fr.tex
\oldpage[III]{203}
\begin{remarks}[7.7.9]
\textnormal{(i)} Dans (7.7.6), (7.7.7), (7.7.8), la formation des
$\mathscr O_Y$-Modules $\mathscr Q$, $\mathscr R$, $\mathscr N$ \emph{permute aux
changements de base}. Par exemple (gardant les notations de (7.7.2)), dans le
cas (7.7.6), on a, pour tout $\mathscr O_{Y'}$-Module quasi-cohérent
$\mathscr M'$, l'isomorphisme
\[
  f'_*(\mathscr F\otimes_{\mathscr O_Y}\mathscr M')
  \xrightarrow{\sim}
  \mathscr{H}\!om_{\mathscr O_{Y'}}(g^*(\mathscr Q),\mathscr M')
\]
car en vertu de la remarque faite dans (7.7.2), tout revient à voir que l'on a
\[
  \operatorname{Hom}_{\mathscr O_Y}(\mathscr Q,g_*(\mathscr M'))
  =
  \operatorname{Hom}_{\mathscr O_{Y'}}(g^*(\mathscr Q),\mathscr M')
\]
ce qui n'est autre que ($\mathbf{0}_{\mathrm I}$, 4.4.3.1). De même, quand dans (7.7.7)
on remplace $Y$, $f$, $\mathscr M$, $\mathscr F$, $\mathscr F'$ par $Y'$, $f'$,
$\mathscr M'$, $\mathscr F\otimes_{\mathscr O_X}\mathscr O_{X'}$,
$\mathscr F'\otimes_{\mathscr O_X}\mathscr O_{X'}$, il faut remplacer
$\mathscr R$ par $g^*(\mathscr R)$. Enfin, dans (7.7.8), lorsqu'on remplace $X$,
$Y$, $f$, $\mathscr M$, $\mathscr F$ par $X'$, $Y'$, $f'$, $\mathscr M'$,
$\mathscr F\otimes_{\mathscr O_X}\mathscr O_{X'}$, et $\mathscr G$ par
$\mathscr G'=\mathscr G\otimes_{\mathscr O_Y}\mathscr O_{X'}$, il faut remplacer
$\mathscr N$ par $g^*(\mathscr N)$ : cela résulte de ce que l'on a encore une
suite exacte
$\mathscr E'_1\longrightarrow\mathscr E'_0\longrightarrow\mathscr G'\longrightarrow0$ avec
$\mathscr E'_i=\mathscr E_i\otimes_{\mathscr O_Y}\mathscr O_{X'}$
et de ce que
$  \check{\mathscr E}'_i
  \otimes_{\mathscr O_{X'}}
  (\mathscr F\otimes_{\mathscr O_Y}\mathscr M')
  =
  \check{\mathscr E}_i
  \otimes_{\mathscr O_X}
  (\mathscr F\otimes_{\mathscr O_Y}\mathscr M)$ ($i=0,1$).

\medskip
\textnormal{(ii)} La condition (ii) de l'énoncé de (7.7.8) relative à
$\mathscr G$ est toujours satisfaite pour un $\mathscr O_X$-Module cohérent
$\mathscr G$ quelconque lorsqu'il existe un $\mathscr O_X$-Module inversible
$Y$-\emph{ample}, par exemple lorsque $Y$ est affine et $f:X\to Y$ est un morphisme
\emph{projectif}. Il suffit alors de noter (compte tenu de (\textbf{II}, 5.5.1)) qu'il existe un
$\mathscr O_X$-Module localement libre de type fini $\mathscr E_0$ tel que
$\mathscr G$ soit isomorphe à un quotient de $\mathscr E_0$ (\textbf{II}, 2.7.10) ; comme
$\mathscr E_0$ et $\mathscr G$ sont cohérents, il en est de même du noyau
$\mathscr G_1$ de $\mathscr E_0\to\mathscr G$, et en appliquant le même résultat
à $\mathscr G_1$, on obtient bien une suite exacte
$\mathscr E_1\to\mathscr E_0\to\mathscr G\to0$ où $\mathscr E_0$ et
$\mathscr E_1$ sont localement libres de type fini.

\medskip
\textnormal{(iii)} Nous prouverons au chap. V que, dans (7.7.8), l'hypothèse
restrictive (ii) est \emph{surabondante}.
\end{remarks}

\begin{proposition}[7.7.10]
\textnormal{(critères locaux pour la propriété d'échange).} Sous les conditions
générales de \textup{(7.7.5)}, soient $y$ un point de $Y$, $p$ un entier. Les propriétés
suivantes sont équivalentes :
\begin{enumerate}
  \item[\textup{a)}] Le foncteur $T_p^{\mathscr O_y}$ est exact à droite.

  \item[\textup{b)}] L'homomorphisme canonique
    $      T_p^{\mathscr O_y}(\mathscr O_y)
      \longrightarrow
      T_p^{\mathscr O_y}(\boldsymbol{k}(y))$
    est surjectif.

  \item[\textup{c)}] Pour tout entier $n$, l'homomorphisme canonique
    $      T_p^{\mathscr O_y}
      (\mathscr O_y/\mathfrak m_y^{n+1})
      \longrightarrow
      T_p^{\mathscr O_y}(\boldsymbol{k}(y))$
    est surjectif.
\end{enumerate}

De plus, l'ensemble des $y\in Y$ vérifiant ces conditions est le plus grand
ensemble ouvert $U$ de $Y$ tel que $\mathscr T_p^U$ soit exact à droite.
\end{proposition}

\begin{proof}
Compte tenu de (7.7.4), l'équivalence de \emph{a)}, \emph{b)} et \emph{c)} résulte de (7.4.7) et
(7.5.2). Le fait que l'ensemble $U$ où $T_p^{\mathscr O_y}$ est exact à droite
soit ouvert est aussi conséquence de (7.4.4), et inversement si
$\mathscr T_p^V$ est exact à droite, il en est de même de
$T_p^{\mathscr O_y}$ pour tout $y\in V$, par la condition \emph{c)} de (7.7.5) et
(7.6.1).
\end{proof}

\begin{corollary}[7.7.11]
Si $\mathscr T_p$ est exact à droite \textup{(resp.} à gauche\textup{)}, alors, pour tout morphisme
$g:Y'\to Y$, $\mathscr T_p^{Y'}$ est exact à droite \textup{(resp.} à gauche\textup{)}. La
réciproque est vraie lorsque le morphisme $g$ est fidèlement plat.
\end{corollary}

